\documentclass{article}
\usepackage[T1]{fontenc}
\usepackage{lmodern}
\usepackage{graphicx}
\usepackage{amsmath,amssymb}
\usepackage[a4paper,margin=2cm]{geometry}
\usepackage[absolute,overlay]{textpos}
\setlength{\TPHorizModule}{1mm}
\setlength{\TPVertModule}{1mm}
\usepackage[indonesian]{babel}
\usepackage{hyperref}
\setlength{\emergencystretch}{2em}
% unit: Halaman 68

\begin{document}

% === Halaman 68 (llm) ===

2. Selanjutnya, ketiga register didetak selama 22 putaran tambahan. Selama 22 putaran tersebut, 22-bit dari nomor $frame$ ($F_n$) dicampur dengan bit register berdasarkan skema berikut:

Pada putaran $0 \le i < 22$,

\begin{itemize}
    \item bit $F_n[i]$ ditambahkan ke bit $LSB$ dari setiap register $R$ dengan menggunakan operasi XOR:
    \[ R[0] = R[0] \oplus F_n[i] \]
    \item tiap register kemudian didetak
\end{itemize}

Isi register pada akhir putaran menyatakan kondisi awal untuk pembangkitan $frame$ sepanjang 228-bit.

\end{document}
